\section{The exact finite arithmetic certificate at thirteen}
\label{sec-f13ds-formal-scope}

The module
\texttt{F13Fixed\allowbreak{}Curve\allowbreak{}Arithmetic.lean}
contains fifteen new theorems about literal arithmetic over
$\mathbb F_{13}$. These are distinct from the twenty-four selector
theorems in the preceding publication. The new module imports no old
ABC project module. Its finite data model, polynomial identities and
chart predicates are described here in full; the geometric and
global interpretations remain separate ordinary proofs.

\paragraph{An identity of polynomials.}
Let $F=\mathbb Z/13\mathbb Z$ and
\[
 S(X)=X^6-27X^4+99X^2-9\in F[X].
\]
The implementation separately defines its literal evaluation $S(z)$
and proves agreement with polynomial evaluation for every $z\in F$.
It also proves the equality
\begin{equation}
 (10+8X^2)S(X)+(6X+3X^3)S'(X)=1
 \qquad\hbox{in }F[X].
 \label{eq-f13ds-bezout}
\end{equation}
This is a coefficient identity in
\texttt{Polynomial (ZMod 13)}, not an assertion of equality at only
thirteen field elements. For an explicit verification, expansion
over the integers gives
\[
 \begin{split}
 &(10+8X^2)(X^6-27X^4+99X^2-9)\\
 &\quad +(6X+3X^3)(6X^5-108X^3+198X)-1\\
 &=13(-7+162X^2+36X^4-38X^6+2X^8).
 \end{split}
\]
The Lean proof differentiates the literal polynomial and uses this
linear combination of the coefficient equality $13=0$ in $F[X]$.

\paragraph{Complete finite sets and their cardinalities.}
The full affine set satisfying $W^2=S(z)$ is exactly
\[
 \begin{array}{c|c}
 z&\text{permitted }W\\ \hline
 0&2,11\\
 1,12&5,8\\
 3,6,7,10&3,10
 \end{array}
\]
and has fourteen elements. The two affine equations
\[
 \mathcal E_{a,b}(x,y):\quad y^2=x^3+ax+b,
 \qquad (a,b)=(-9,-9),\;(-189,999)
\]
also each have fourteen solutions in $F^2$.
These are affine cardinalities; the module does not construct the
elliptic group or prove its order from these counts.

For the complete finite chart model define the disjoint sum
\[
 \mathcal M=(F\times F)\sqcup F.
\]
An element in its first summand is valid when $W^2=S(z)$;
an element $\varepsilon$ in its second summand is valid when
$\varepsilon^2=1$. The latter condition is proved equivalent to
$\varepsilon=1$ or $\varepsilon=-1$. The valid elements of
$\mathcal M$ therefore number sixteen, as proved by a literal
finite-set count. This disjoint-sum definition is not presented
as a formal construction of a smooth projective curve.

\paragraph{Actual affine coordinates on all three charts.}
The field functions used in the source are
\[
 \begin{aligned}
 x_1(z)&=\frac{z^2-9}{4},&
 y_1(W)&=\frac W8,\\
 x_2(z)&=\frac{33-9/z^2}{4},&
 y_2(z,W)&=-\frac{9W}{8z^3}.
 \end{aligned}
\]
For every $z,W\in F$ with $W^2=S(z)$ and $z\ne0$, the source proves
both actual curve equations
\[
 \mathcal E_{-9,-9}(x_1(z),y_1(W)),\qquad
 \mathcal E_{-189,999}(x_2(z),y_2(z,W)).
\]
For the zero chart it proves, for every $W^2=S(0)$,
\[
 \mathcal E_{-9,-9}\left(-\frac94,\frac W8\right).
\]
For every infinity sign $\varepsilon^2=1$, it proves
\[
 \mathcal E_{-189,999}\left(\frac{33}{4},-\frac{9\varepsilon}{8}\right).
\]
The missing affine coordinates in the last two cases are handled
by the explicit finite predicates below. The geometric statement
that the other projection is the elliptic point at infinity is
not a theorem about an abstract elliptic group in this module.

\paragraph{Literal division numerators and the finite obstruction.}
For $a,b,x\in F$ define
\[
 \begin{aligned}
 \Psi_{a,b}(x)&=3x^4+6ax^2+12bx-a^2,\\
 D_{a,b}(x)&=x^6+5ax^4+20bx^3-5a^2x^2\\
           &\hspace{9mm}-4abx-8b^2-a^3,\\
 \Phi_{a,b}(x)&=x\Psi_{a,b}(x)^2
                  -8(x^3+ax+b)D_{a,b}(x).
 \end{aligned}
\]
For arbitrary $a,b,x,y\in F$ satisfying the actual equation
$y^2=x^3+ax+b$, the formal substitution theorem is
\[
 \Phi_{a,b}(x)
 =x\Psi_{a,b}(x)^2-(4yD_{a,b}(x))(2y).
\]
Put
\[
 N_2(x)=\Phi_{-189,999}(x)-11\Psi_{-189,999}(x)^2.
\]
The complete nonzero-chart exclusion is
\[
 \begin{gathered}
 z,W\in F,\quad W^2=S(z),\quad z\ne0,\quad
 \Psi_{-9,-9}(x_1(z))=0\\
 \Longrightarrow\quad N_2(x_2(z))\ne0.
 \end{gathered}
\]
The remaining two exact values are
\[
 \Psi_{-9,-9}(-9/4)=7,\qquad N_2(33/4)=9.
\]
They are identities in $F$ and in particular are nonzero.

For clarity, the final assembled predicates on $\mathcal M$ are
\[
 \begin{array}{c|c|c}
 p&\mathsf{FirstZero}(p)&\mathsf{SecondTarget}(p)\\ \hline
 (z,W)&\Psi_{-9,-9}(x_1(z))=0&
             z\ne0\ \text{and}\ N_2(x_2(z))=0\\
 \varepsilon&\text{true}&N_2(33/4)=0 .
 \end{array}
\]
The last theorem proves for every element of the finite model
\[
 \mathsf{Valid}(p)\ \Longrightarrow\
 \neg\bigl(\mathsf{FirstZero}(p)\ \text{and}\
                 \mathsf{SecondTarget}(p)\bigr).
\]
In particular the second condition on an affine element with $z=0$
is false. The chart implications and this assembled theorem use
all finite inputs, not a sample of rational points. The ordinary
division-polynomial interpretation connects these numerators with
the two signs of the elliptic target of abscissa eleven; that
interpretation as multiplication by three remains outside the
formal statements.

\paragraph{The exact fifteen declarations.}
The preceding statements correspond to the following theorem names,
each with its own explicit axiom query:
\begin{center}
\small
\begin{tabular}{r l}
1&\texttt{sextic\_eval}\\
2&\texttt{sextic\_polynomial\_bezout}\\
3&\texttt{affine\_table\_exact}\\
4&\texttt{affine\_count}\\
5&\texttt{elliptic\_affine\_counts}\\
6&\texttt{infinity\_signs}\\
7&\texttt{complete\_finite\_model\_count}\\
8&\texttt{nonzero\_projection\_equations}\\
9&\texttt{zero\_projection\_equation}\\
10&\texttt{infinity\_projection\_equation}\\
11&\texttt{numerator\_substitution}\\
12&\texttt{nonzero\_chart\_exclusion}\\
13&\texttt{zero\_chart\_first\_psi}\\
14&\texttt{infinity\_chart\_second\_numerator}\\
15&\texttt{complete\_finite\_exclusion}
\end{tabular}
\end{center}

\paragraph{Actual verification and its boundary.}
The author compiled a literal source copy in a fresh isolated
project, using Lean 4.32.0 and the pinned Mathlib cache. Root
independently repeated that fresh compilation. Both runs passed
with warnings treated as errors, and all fifteen printed axiom
queries contain only \texttt{propext},
\texttt{Classical.choice} and \texttt{Quot.sound}.
The critical-bottleneck reviewer separately read the complete
source, ordinary scope, verifier, manifest and compiler log and
recomputed their hashes. That source review is not counted as
another compiler execution.

The source SHA256, with the following two lines concatenated, is
\begin{center}
\small\ttfamily
a459a5f16e8f0b8f6a2dd2dbfd742c67\\
2721f3963220463eed10e8b456e006ec
\end{center}
The two successful per-run records are identified by
\begin{center}
\small\ttfamily
20260908T055706020261Z\\
\textrm{and }20260908T060020018236Z
\end{center}
in the nineteenth-round
\texttt{next\_verification} directory. They bind source, prior
ordinary scope, verifier, full log, toolchain, dependency pin,
all queries and the fresh olean. Earlier failed attempts remain
separate failed records and are not proofs.

Finite division expressions require direct kernel reduction through
the \texttt{ZMod} inverse. The source uses
\texttt{decide +kernel}; it does not use native evaluation,
\texttt{sorry}, or a new axiom. The polynomial proof uses algebraic
tactics on the actual coefficients.

These fifteen theorems do not formalize the projective curve,
abstract elliptic multiplication by three, good reduction, the
global prime-to-five index, ranks, the two-descent, $p$-adic heights
or zero certificates, or rational-point completeness. Those are
separately stated ordinary mathematical inputs and conclusions.
No such bridge is hidden as a hypothesis or axiom in this module.
Neither this finite certificate nor the previously published
twenty-four selector theorems prove the ABC conjecture.
